\documentclass[11pt,a4paper]{article}
\usepackage{fontspec}
\usepackage[margin=2.5cm]{geometry}
\usepackage{hyperref}
\usepackage{parskip}
\usepackage{titlesec}
\usepackage{enumitem}
\hypersetup{colorlinks=true,urlcolor=blue}
\titleformat{\section}{\large\bfseries}{}{0pt}{}
\titlespacing*{\section}{0pt}{1.4em}{0.3em}
\pagestyle{empty}

\begin{document}

{\LARGE\bfseries CS491 -- Fall 2026 -- Project Information Form\par}
\vspace{0.4em}
Project information submitted for the CS491 Design project.

\section*{Name of the Project}
GENIE5 (GEne Network Inference from Evidence)

\section*{One paragraph description of the project}
Genie5 turns a biomedical research question into a directed, evidence-backed map of gene regulation. Statistical analysis can reveal that two genes are linked, but rarely which one regulates the other, so researchers still settle the question by hand-searching the literature, one gene pair at a time. Genie5 automates that review. It collects evidence from scientific papers and curated biological databases, reconciles conflicting reports, and determines the direction and type of each interaction, such as activation or repression. Researchers receive an interactive 2D/3D network that they can explore, verify and export to their existing tools, including CRNA for analysis on their own expression data.

\section*{Supervisor}
Prof. Can Alkan, Department of Computer Engineering, Bilkent University

\section*{Innovation Expert}
Asst. Prof. Gurkan Bebek, Case Center for Proteomics and Bioinformatics, Case Western Reserve University

\section*{Project Web Page}
\url{https://twin-pathways.github.io/}

\section*{Team Members}
\begin{tabular}{@{}ll@{}}
\textbf{Student ID} & \textbf{Name} \\[0.3em]
22302668 & Sırrı Kaan Oğuzkan \\
22302229 & Derin Ceylin Atalar \\
22302289 & Ayça Topkaya \\
22303934 & Alper Koçyiğit \\
22302218 & Ömer Turab Memiş \\
\end{tabular}

\end{document}
